% Libro 4 · anexos: guía rápida, plan de estudio y simulacro final
\makeatletter
\cleardoublepage\renewcommand{\zuka@mundoactual}{Guía rápida}
\phantomsection\addcontentsline{toc}{chapter}{Guía rápida}
{\sffamily\fontsize{26}{30}\selectfont\firalight Guía rápida\par}\vspace{2pt}
{\color{zlime}\rule{18mm}{1.6pt}}\par\vspace{6pt}

\begin{multicols}{2}\small
\kicker{Inferencia (1)}
\begin{itemize}
  \item Conclusión con pruebas; la mínima gana.
  \item Presupuesto: sobrevive a la negación.
  \item Cadena causal: directa u original.
  \item Actitudes: los hechos delatan; intensidad justa.
\end{itemize}
\kicker{Aplicación (2)}
\begin{itemize}
  \item Extrae el principio «si X, entonces Y».
  \item Analogía: nombra la relación con una frase.
\end{itemize}
\kicker{Voces (3)}
\begin{itemize}
  \item ¿Para qué la cita? Apoyar, oponer, ilustrar.
  \item Fuente: quién, qué interés, cómo lo sabe.
\end{itemize}
\columnbreak
\kicker{Argumentación (4)}
\begin{itemize}
  \item Árbol: tesis, argumentos, objeción, refutación.
  \item Intención: qué quiere lograr. Tono: adjetivos y adverbios.
  \item Hecho: verificable. Opinión: valorable.
\end{itemize}
\kicker{Figuras (5)}
\begin{itemize}
  \item Comparación (como), metáfora, personificación, hipérbole, antítesis, metonimia, anáfora.
  \item Ironía: choque entre lo dicho y la situación.
\end{itemize}
\kicker{Dos textos (6)}
\begin{itemize}
  \item Tabla de dos columnas: tema, tesis, argumentos, tono.
  \item ¿Qué diría A? Usa la lógica de A, no la tuya.
\end{itemize}
\end{multicols}

\clearpage\renewcommand{\zuka@mundoactual}{Plan de estudio}
\phantomsection\addcontentsline{toc}{chapter}{Plan de estudio con la app}
{\sffamily\fontsize{26}{30}\selectfont\firalight Plan de seis semanas\par}\vspace{2pt}
{\color{zlime}\rule{18mm}{1.6pt}}\par\vspace{6pt}
{\small Como en el Libro 3: una lección diaria del libro con su lección en Zuka, y la evaluación del mundo los sábados. En
la semana 6, el simulacro final reproduce la tarea más exigente del examen: dos textos que dialogan.\par}
\vspace{6pt}
\begin{center}\small
\begin{tabularx}{\linewidth}{@{}>{\sffamily\bfseries\color{zolive}}lX>{\raggedleft\arraybackslash}p{24mm}@{}}
\toprule
\textbf{Semana} & \textbf{Lecciones} & \textbf{Evaluación}\\\midrule
1 & 1.1 Implícitas · 1.2 Presupuestos · 1.3 Causas · 1.4 Actitudes & Mundo 1\\
2 & 2.1 Aplicar · 2.2 Analogías · 3.1 Citas · 3.2 Fuentes & Mundos 2 y 3\\
3 & 4.1 Tesis · 4.2 Intención y tono · 4.3 Hechos y opiniones & Mundo 4\\
4 & 5.1 Figuras · 5.2 Ironía & Mundo 5\\
5 & 6.1 Coincidencias · 6.2 Qué diría el otro & Mundo 6\\
6 & Simulacro final · repaso de errores de los libros 3 y 4 & Simulacro\\
\bottomrule
\end{tabularx}
\end{center}

\clearpage\renewcommand{\zuka@mundoactual}{Simulacro final}
\setcounter{ej}{0}
\phantomsection\addcontentsline{toc}{chapter}{Simulacro final}
\immediate\write\zuka@sol{\noexpand\solmun{Simulacro final}}
\cabecera{SIMULACRO FINAL · LIBRO 4}{Leer lo que quiere decir}{dos textos · 10 preguntas\\18 minutos}\vspace{4pt}
\begin{texto}{Texto 1 · Elogio del aburrimiento}{Ensayo · texto original escrito para este libro}
\pa Nunca habíamos tenido tan pocos momentos de aburrimiento. En la fila del banco, en el bus, en la sala de espera, basta
sacar el teléfono para llenar cualquier vacío. Parece un avance. No estoy seguro de que lo sea.

\pa El aburrimiento es incómodo, pero es fértil. Cuando no hay nada que mirar, la mente empieza a divagar, a combinar
recuerdos, a imaginar. Muchas ideas nacen en esos ratos muertos que hoy hemos aprendido a evitar a toda costa.

\pa Los niños lo saben mejor que nadie: después de quejarse de que «no hay nada que hacer», inventan un juego. Les
quitamos esa oportunidad cada vez que les damos una pantalla para que no se aburran.

\pa Propongo, entonces, algo modesto: defender algunos minutos de aburrimiento al día. No como castigo, sino como
espacio para que aparezca lo que no sabíamos que estaba ahí.
\end{texto}
\begin{texto}{Texto 2 · El aburrimiento no es un lujo para todos}{Columna · texto original escrito para este libro}
\pa Está de moda elogiar el aburrimiento. Quienes lo hacen suelen tener tiempo libre, una casa tranquila y la seguridad de
que, después de aburrirse, nada malo les pasará.

\pa Para muchas personas, en cambio, los «ratos muertos» son pocos y están ocupados por preocupaciones: cómo pagar el
arriendo, cómo llegar a tiempo al segundo trabajo. Para ellas, el teléfono en el bus no es una evasión de la creatividad,
sino a veces el único descanso del día.

\pa No niego que la mente necesita pausas. Pero antes de pedir a todos que se aburran más, habría que preguntarse quién
tiene el privilegio de hacerlo.
\end{texto}
\begin{ejercicios}
\item La tesis del texto 1 es que
\opcionesl{los teléfonos deben prohibirse.}{el aburrimiento, aunque incómodo, favorece la creatividad y conviene defenderlo.}{los niños se aburren demasiado.}{las filas son aburridas.}
\sol{B. (4.1)}
\item En el texto 1, el ejemplo de los niños \pf{3} sirve para
\opcionesl{ilustrar que el aburrimiento lleva a inventar.}{criticar a los niños.}{presentar una objeción.}{dar un dato estadístico.}
\sol{A. (3.1)}
\item La frase «no estoy seguro de que lo sea» (texto 1, \S\,1) presupone que
\opcionesl{muchos consideran un avance poder evitar el aburrimiento.}{el autor odia los teléfonos.}{nadie se aburre.}{el banco tiene filas largas.}
\sol{A. (1.2)}
\item Las comillas en «ratos muertos» (texto 2, \S\,2) indican que el autor
\opcionesl{cita una expresión asociada a quienes elogian el aburrimiento y la pone en duda.}{inventa un término científico.}{cita a un médico.}{usa una metáfora propia.}
\sol{A. (5.2)}
\item El tono del texto 2 frente al elogio del aburrimiento es
\opcionesd{crítico}{entusiasta}{indiferente}{humorístico}
\sol{A. (4.2)}
\item ¿En qué coinciden ambos textos?
\opcionesl{En que la mente necesita pausas.}{En que el teléfono es el único descanso.}{En que el aburrimiento es un privilegio.}{En que hay que prohibir las pantallas a los niños.}
\sol{A. El texto 2: «No niego que la mente necesita pausas». (6.1)}
\item El autor del texto 2 respondería a la propuesta final del texto 1 que
\opcionesl{es perfecta para todos.}{no todos tienen las condiciones para dedicar ese tiempo al aburrimiento.}{el aburrimiento es dañino.}{los niños deben trabajar.}
\sol{B. (6.2)}
\item ¿Cuál situación ilustra mejor la idea del texto 1?
\opcionesl{Una persona revisa el teléfono en cada semáforo.}{En un viaje largo sin señal, una estudiante empieza a escribir una historia.}{Un trabajador usa el bus para dormir.}{Un niño ve dibujos animados.}
\sol{B. (2.1)}
\item «Les quitamos esa oportunidad» (texto 1, \S\,3): «esa oportunidad» es
\opcionesl{la de tener una pantalla.}{la de aburrirse y terminar inventando un juego.}{la de quejarse.}{la de ir al banco.}
\sol{B. Referencia. (Libro 3, 5.1)}
\item ¿Cuál es una opinión y no un hecho en el texto 2?
\opcionesl{Hay personas con dos trabajos.}{Está de moda elogiar el aburrimiento.}{Habría que preguntarse quién tiene el privilegio de aburrirse.}{Algunas personas usan el teléfono en el bus.}
\sol{C. Es una valoración y una recomendación. (4.3)}
\end{ejercicios}
\makeatother
